\chapter{Spring 核心：容器、Bean、依赖注入、配置}

这一关学 Spring 最核心的概念。后面几关的每一个类，都是按这一关的方式交给 Spring 管理的。

\partinfo{%
  \code{config/AppProperties.java}：2.1、2.2\\
  \code{config/AppConfig.java}：2.1、2.3\\
  \code{station/domain/StationFinder.java}：2.4%
}{\code{./mvnw test -Dtest='Part2*'}}

\section{知识点}

\subsection{Spring 容器和 Bean}

在普通 Java 程序里，你自己 \code{new} 对象，自己把对象传给需要它的地方：

\begin{java}
BookRepository repo = new BookRepository(dataSource);
BookService service = new BookService(repo);
BookController controller = new BookController(service);
\end{java}

项目一大，这种“谁 new 谁、谁传给谁”的代码会非常多。Spring 的做法是：\textbf{你只声明有哪些类、每个类需要什么，Spring 负责创建对象并把它们连起来。}

\begin{itemize}
  \item 由 Spring 创建和管理的对象叫 \textbf{Bean}。
  \item 管理所有 Bean 的东西叫\textbf{容器}（ApplicationContext）。
  \item 默认每个 Bean 只创建一个实例（单例），所有用到它的地方共享同一个。
\end{itemize}

\subsection{怎么把一个类注册成 Bean}

\textbf{方法一：在类上加注解，让 Spring 扫描到。}

\begin{java}
@Component
public class PriceCalculator { ... }
\end{java}

启动类 \code{DublinBikesSpringApplication} 上的 \code{@SpringBootApplication} 会扫描它所在的包（\code{ie.dublinbikes}）以及所有子包，找到带这类注解的类，自动创建 Bean。

\code{@Component} 是通用的。还有几个意思更具体的，效果一样，只是让代码更好读：

\begin{center}
\begin{tabular}{ll}
\toprule
注解 & 用在 \\
\midrule
\code{@Component} & 通用组件 \\
\code{@Service} & 业务逻辑层 \\
\code{@Repository} & 数据访问层 \\
\code{@RestController} & 提供 REST 接口的类（Part 6） \\
\bottomrule
\end{tabular}
\end{center}

\textbf{方法二：在配置类里写一个方法，返回值就是 Bean。}有些类不是你写的（比如 JDK 的 \code{Clock}），没办法在上面加注解，就用这种方式：

\begin{java}
@Configuration                 // 表示这是一个配置类
public class MyConfig {

    @Bean                      // 这个方法的返回值会注册成 Bean
    public Random random() {
        return new Random(42);
    }
}
\end{java}

\subsection{依赖注入（DI）：通过构造函数}

一个 Bean 需要用到另一个 Bean 时，\textbf{在构造函数参数里声明就行}，Spring 创建它的时候会自动把需要的 Bean 传进来：

\begin{java}
@Service
public class BookService {

    private final BookRepository repository;

    public BookService(BookRepository repository) {   // Spring 自动传入 BookRepository 的 Bean
        this.repository = repository;
    }
}
\end{java}

只有一个构造函数时，不需要加任何注解。这种写法叫\textbf{构造函数注入}，是官方推荐的方式。你会在网上看到另一种写法 \code{@Autowired private BookRepository repository;}（字段注入），不推荐，原因见本章最后的“想一想”。

\subsection{配置文件绑定到 record}

\code{src/main/resources/application.yml} 是应用的配置文件。打开看一下，里面有这样一段：

\begin{yamlcode}
app:
  jcdecaux:
    api-key: ${JCDECAUX_API_KEY:}
  scraper:
    enabled: ${SCRAPER_ENABLED:false}
    interval: 5m
\end{yamlcode}

\code{${JCDECAUX_API_KEY:}} 的意思是：读取环境变量 \code{JCDECAUX_API_KEY}，没有的话用冒号后面的默认值（这里是空字符串）。

在代码里读这些配置，最好的方式是把它们\textbf{绑定到一个 record}：

\begin{java}
@ConfigurationProperties(prefix = "library")
public record LibraryProperties(String name, int maxLoans, Duration loanPeriod) {}
\end{java}

\begin{yamlcode}
library:
  name: Central
  max-loans: 5          # 配置里用 kebab-case，自动对应 Java 的 maxLoans
  loan-period: 14d      # 字符串自动转成 Duration（支持 30s、5m、2h、14d）
\end{yamlcode}

嵌套的配置对应嵌套的 record：\code{app.jcdecaux.api-key} 对应 \code{AppProperties} 里 \code{Jcdecaux} record 的 \code{apiKey}。

光写 \code{@ConfigurationProperties} 还不够，还要告诉 Spring “请把这个 record 注册成 Bean”。在某个 \code{@Configuration} 类上加：

\begin{java}
@EnableConfigurationProperties(LibraryProperties.class)
\end{java}

之后任何 Bean 都可以在构造函数里要一个 \code{LibraryProperties}。

\textbf{默认值}：配置里没写的字段会是 null（基本类型是 0 或 false）。要给默认值，在 record 参数上加 \code{@DefaultValue}（包名 \code{org.springframework.boot.context.properties.bind}）：

\begin{java}
public record LibraryProperties(
        @DefaultValue("Central") String name,
        @DefaultValue("5") int maxLoans) {}
\end{java}

嵌套 record 有个坑：如果配置文件里完全没有 \code{library.xxx} 这一段，嵌套的 record 本身会是 null。在\textbf{外层}的参数上加一个不带值的 \code{@DefaultValue}，就能让 Spring 用默认值创建它：

\begin{java}
public record Outer(@DefaultValue Inner inner) {}
\end{java}

\subsection{测试是怎么验证的}

打开 \code{src/test/java/ie/dublinbikes/config/Part2SpringCoreTest.java} 看一下。\code{ApplicationContextRunner} 会启动一个只包含 \code{AppConfig} 的小容器，用 \code{withPropertyValues(...)} 模拟配置文件，然后检查容器里有没有你期望的 Bean、值对不对。

\section{题目}

\begin{task}{2.1\quad 配置绑定}
让 \code{application.yml} 里 \code{app.*} 的配置能绑定到 \code{AppProperties}：
\begin{itemize}
  \item \code{AppProperties} 对应前缀 \code{app}。
  \item \code{AppConfig} 是一个配置类，并且启用 \code{AppProperties} 的绑定。
\end{itemize}

\begin{example}
\begin{plaincode}
配置：app.jcdecaux.api-key=test-key
      app.jcdecaux.contract=lyon
      app.scraper.interval=30s

结果：props.jcdecaux().apiKey()   == "test-key"
      props.jcdecaux().contract() == "lyon"
      props.scraper().interval()  == Duration.ofSeconds(30)
\end{plaincode}
\end{example}
\end{task}

\begin{task}{2.2\quad 默认值}
什么都没配置时：

\medskip
\begin{tabular}{ll}
\toprule
字段 & 默认值 \\
\midrule
\code{jcdecaux().apiKey()} & \code{null}（不需要默认值） \\
\code{jcdecaux().contract()} & \code{"dublin"} \\
\code{jcdecaux().baseUrl()} & \code{"https://api.jcdecaux.com/vls/v1"} \\
\code{scraper().enabled()} & \code{false} \\
\code{scraper().interval()} & 5 分钟 \\
\bottomrule
\end{tabular}

\medskip
而且 \code{jcdecaux()} 和 \code{scraper()} 本身不能是 null。
\end{task}

\begin{task}{2.3\quad Clock Bean}
在 \code{AppConfig} 里把 \code{clock()} 方法的返回值注册成 Bean，返回一个使用 \code{Europe/Dublin} 时区的系统时钟（看一下 \code{Clock.system(...)} 和 \code{ZoneId.of(...)}）。

为什么要把时钟做成 Bean，而不是在代码里直接用 \code{Instant.now()}？见本章最后的“想一想”。
\end{task}

\begin{task}{2.4\quad 组件扫描}
让 Spring 能通过组件扫描自动发现 \code{StationFinder}。Part 6 的 \code{StationQueryService} 会在构造函数里要一个 \code{StationFinder}。
\end{task}

\begin{think}
\begin{enumerate}
  \item 为什么推荐构造函数注入，不推荐 \code{@Autowired} 字段注入？（提示：\code{final} 字段、不启动 Spring 能不能写单元测试、循环依赖）
  \item Bean 默认是单例。如果一个单例 Bean 里有一个会被修改的字段，多个请求同时访问会出什么问题？
  \item 代码里直接调用 \code{Instant.now()}，和注入一个 \code{Clock} 再调用 \code{Instant.now(clock)}，测试的时候有什么区别？
  \item \code{@Component} 和 \code{@Bean} 分别适合什么场景？
\end{enumerate}
\end{think}
